\section{Review Questions}

\begin{enumerate}
\item Define an inverse matrix. Why are both identities $AA^{-1}=I$ and $A^{-1}A=I$ relevant in a noncommutative setting?

\item Explain an inverse as the reversal of a matrix action. What kind of information loss prevents such a reversal?

\item State the invertibility criterion involving the determinant and interpret both cases $\det A\neq0$ and $\det A=0$ structurally.

\item Derive the inverse formula for
\[
A=\begin{pmatrix}a&b\\c&d\end{pmatrix}
\]
when $ad-bc\neq0$. Identify exactly where the determinant enters and why division by zero is impossible.

\item Use the $2\times2$ formula to invert one simple matrix and verify the answer by multiplication. What does the verification test?

\item Explain the Gauss--Jordan construction
\[
[A\mid I]\longrightarrow[I\mid A^{-1}].
\]
Why must every row operation be applied to both blocks?

\item What does failure to reduce the left block to the identity reveal? Describe the connection among missing pivots, a zero determinant, and non-invertibility.

\item Define the cofactor matrix and the adjugate. For a general $2\times2$ matrix, show how the adjugate recovers the familiar inverse formula.

\item Compare the $2\times2$ formula, Gauss--Jordan elimination, and the adjugate formula:
\begin{enumerate}
\item which is the simplest in dimension two;
\item which is the principal general hand method;
\item which best exposes the determinant--cofactor structure?
\end{enumerate}

\item Solve a matrix equation such as $AXB=C$ under suitable invertibility assumptions. Explain the order in which inverses must be applied and why ordinary scalar-style cancellation would be unsafe.
\end{enumerate}
